% Pista ana-23j-q5 · Anàlisi
% Procedència: PAU Catalunya 2023 · Convocatòria de juny · Sèrie 1 · Qüestió 5
\documentclass[11pt,a4paper]{article}
\usepackage{pista}
\pistainfo{Catalunya 2023}{Convocatòria de juny}{Sèrie 1}

\begin{document}

\section*{\textcolor{blauPista}{Qüestió 5}}

\noindent La Núria té un jardí rectangular i vol fer-hi un tancat (rectangular o quadrat) de $8\,\text{m}^{2}$ per al seu gos. Ha pensat de posar el tancat tocant al mur del jardí, tal com es mostra a la figura de la dreta, per estalviar-se així un dels quatre costats. El preu de la tanca que vol fer servir és de $2{,}5\,\euro/\text{m}$.

\subsection*{\textit{a)} \normalfont Quines dimensions ha de tenir el tancat perquè el cost sigui mínim? Quin és aquest cost mínim? \hfill\textnormal{[1,75~punts]}}

\pista{Identifica clarament les dues dimensions del tancat: una paral·lela al mur (per exemple $x$) i una perpendicular ($y$). La condició de l'àrea, $x \cdot y = 8$, et permet expressar una variable en funció de l'altra. A continuació, escriu la longitud de tanca que cal instal·lar: ves amb compte, perquè només són TRES costats (el quart costat el fa el mateix mur del jardí). Multiplica per $2{,}5\,\euro/\text{m}$ i obtindràs la funció cost $C$ en funció d'una sola variable. A partir d'aquí, és un problema d'optimització estàndard: deriva $C$, resol $C'(x) = 0$, i comprova que el valor obtingut correspon efectivament a un mínim. Recorda que la teva resposta final ha d'incloure les \emph{dues} dimensions del tancat \emph{i} el cost mínim.}

\subsection*{\textit{b)} \normalfont Si manteniu la forma rectangular o quadrada del tancat i feu que un dels vèrtexs del jardí coincideixi amb un vèrtex del tancat, quants euros us podeu estalviar? Raoneu com posaríeu el tancat i justifiqueu amb càlculs matemàtics les dimensions de la vostra proposta. \hfill\textnormal{[0,75~punts]}}

\pista{Si poses un vèrtex del tancat sobre un vèrtex del jardí, ara aprofites DOS murs del jardí (els dos costats consecutius que conflueixen en aquell vèrtex). Per tant, només et caldrà construir DOS costats de tanca, no tres. Refés el plantejament de l'apartat anterior amb aquesta nova configuració: la condició de l'àrea segueix essent la mateixa, $x \cdot y = 8$, però la longitud de tanca és ara $x + y$. Optimitza igual que abans i, finalment, compara el cost mínim obtingut amb el de l'apartat a) per saber l'estalvi en euros.}

\end{document}
